\scanpage{124}
prenant le produit cartésien de
\diagram{\begin{tikzcd}
\{\pm1\}^{A_1}\arrow[r,hook]&\mathbb U^{A_1}
\end{tikzcd}}
\margin{hom. évident déduit de $\{\pm1\}\hookrightarrow\mathbb U$}
\diagram{\begin{tikzcd}
\{\pm1\}\arrow[r,hook]\arrow[dr,"{\text{diag}}"']&\mathbb U^{J_{\mathrm I}}\\
&\{\pm1\}^{J_{\mathrm I}}\arrow[u,hook,"{\text{incl. can}}"']
\end{tikzcd}}
\diagram{\begin{tikzcd}
\{\pm1\}^{J_{\mathrm{III}}}\arrow[r,hook]&\mathbb U^{J_{\mathrm{III}}}
\end{tikzcd}}
\margin{hom. évident.}

Ceci posé, \unsure{enf}, on aura
\[
SP_G^Y\simeq\mathbb U^{\widehat A}\mathop{\wedge}_{\gamma_G}
\underbrace{\varepsilon_{G,Q}}_{\gamma_G\text{-tors.}}
\]

Dans le cas où $J_{\mathrm I}=\varnothing$, on \inserted{peut} poser si on préfère
\[
\varepsilon_{G,Q}=\varepsilon_{G,Y_Q}=\varepsilon_{\widehat A_1}\times\varepsilon_{J_{\mathrm{III}}}
\]
faisant l'économie du facteur $\omega_Q$, \unsure{définissons}
\[
\gamma_G=\{\pm1\}^{\widehat A_1}\times\{\pm1\}^{J_{\mathrm{III}}}
=\{\pm1\}^{\widehat A_1\amalg J_{\mathrm{III}}}
\]
et on \unsure{conservera|conserve} l'expression ci-dessus pour $SP_G^Y$.

Considérons donc
\[
S_0\mathcal T_G^{Y_Q}=S_0\mathcal T_G^Q
\]
\margin{sous-groupoïde plein de $\unsure{\Pi_1|\mathbb{\Pi}_1}(SP_G^Y)$ basé en les éléments de $\varepsilon_{G,Q}$}

Il faudra expliciter ce groupoïde \deleted{fondamental} \inserted{basé sur $\varepsilon_{G,\unsure{Q}}$} par flèches \unsure{génératrices|générateurs} et relations -- j'y reviendrai.

\mspage{123}
